\section{Булевы схемы}

\subsection{Определение}

\begin{defn}{}{}
    \( x_1, x_2, \ldots, x_n \) --- переменные

    \( \lor, \land, \lnot \) --- полная система в \( P_2 \)

    Булева схема --- последовательность \( g_1, \ldots, g_s \in P_2 \)
    \[
        g_i = \left[
            \begin{aligned}
                h_j \lor h_k
                \\
                h_j \land h_K
                \\
                \lnot h_j
            \end{aligned}
        \right.
    \]

    (где \( h_j, h_k \in \{ x_1, \ldots, x_n, g_1, \ldots, g_{i - 1} \)).

    Такая булева схема вычисляет \( g_s \) и ее сложность равна \( s \).
\end{defn}

\begin{example}{}{}
    Схема
    \(
        \lnot x_1, \:
        \lnot x_2, \:
        x_1 \land \lnot x_2, \:
        \lnot x_1 \land x_2, \:
        (x_1 \land \lnot x_2) \lor (\lnot x_1 \land x_2) = x \oplus y
    \)
\end{example}

\begin{example}{}{}
    Схема
    \(
        x_1 \land x_2, \:
        x_1 \lor x_2, \:
        x_3 \land (x_1 \lor x_2), \:
        (x_1 \land x_2) \lor (x_3 \land (x_1 \lor x_2)) = MAJ_3(x_1, x_2, x_3), \:
    \)
\end{example}

\begin{defn}{}{}
    Для \( f : \{ 0, 1 \}^n \to \{ 0, 1 \} \) ее схемная сложность \( c(f) \)
    --- это минимальный размер схемы, вычисляющий \( f \)
\end{defn}

\subsection{Булева схема для сумматора}

Дано:
\( x = \ol{x_1 x_2 \ldots x_n}_2 \),
\( y = \ol{y_1 y_2 \ldots y_n}_2 \)

Цель: вычислить \( x + y \) булевой схемой.

\( z \) --- результат сложения \( x \) и \( y \), \( w \) --- был ли перенос.

\( z_n = x_n \oplus y_n \), \( w_n = x_n \land y_n \)

\( z_{n - 1} = x_{n - 1} \oplus y_{n - 1} \oplus w_n \), \( w_{n - 1} = MAJ(w_n, x_{n - 1}, y_{n - 1} \)

\ldots

\( z_1 = x_1 \oplus y_1 \oplus w_2 \), \( w_1 = MAJ(w_2, x_1, y_1) \)

\( z_0 = w_1 \)

Из примеров: \( c(MAJ_3) \leq 4 \) и \( c(x \oplus y) \leq 5 \).

Поэтому сложность булевой схемы для сумматора \( O(n) \).

\subsection{Общие факты о схемной сложности}

\begin{thrm}{}{}
    Пусть \( f : \{ 0, 1 \}^n \to \{ 0, 1 \} \) --- некоторая функция.
    Тогда \( c(f) = O(n \cdot 2^n) \).
\end{thrm}

При \( f \not\equiv 0 \) можем построить СДНФ.
Для каждого монома \( O(n) \) операций, а таких мономов \( \leq 2^n \) --- доказали.

Для \( f = 0 \) верно \( f = x_1 \land \lnot x_1 \).

\begin{exercise}{}{}
    Доказать, что \( c(f) = O(2^n) \).
\end{exercise}

\begin{exercise}{(довольно сложное)}{}
    Доказать, что \( \displaystyle c(f) = O \left( \frac{2^n}{n} \right) \).
\end{exercise}

\begin{thrm}{}{}
    При \( n \geq 10 \) существует \( f : \{ 0, 1 \}^n \to \{ 0, 1 \} \), для которого \( c(f) \geq \frac{2^n}{10 n} \).
\end{thrm}

\( S = \lfloor \frac{2^n}{10 n} \rfloor \).

Оценим число булевых схем размера \( S \).

\( \lor, \lnot, \land \) можно закодировать двумя битами,
а переменные --- \( \lceil \log_2 (n + S) \rceil \) битами.
Значит на схему мы тратим \( S(2 + \lceil \log_2 (n + S) \rceil) \) бит.
Откуда булевых схем размера \( S \) не более \( 2^{S(2 + \lceil \log_2 (n + S) \rceil)} \).

\(
    \displaystyle \frac{2^n}{10n} > n
\),
поэтому \( S \geq n \)
\begin{gather*}
    S(2 + \lceil \log (n + S) \rceil)
    \leq
    2S + 2S \lceil \log_2 (2S) \rceil
    =
    4S + 2S \lceil \log_2 S \rceil
    \leq
    \\
    \leq
    4S + 2S (\log_2 S + 1)
    =
    6S + 2S \log_2 S
    \leq
    4 S \log_2 S
\end{gather*}

Значит булевых схем размера \( S \) не более \( 2^{4 S \log_2 S} \).

Пусть можем получить все функции схемами размера \( S \):
\begin{gather*}
    2^{2^n} \leq 2^{4 S \log_2 S}
    \\
    2^n \leq 4 S \log_2 S \leq 4 \cdot \frac{2^n}{10 n} \cdot ( n - \log_2 10n )
        <
        4 \cdot \frac{2^n}{10 n} \cdot n
        = 
        \frac{2}{5} \cdot 2^n
        < 2^n
\end{gather*}

Получили противоречие, а значит есть схема сложности более \( S \)

\subsection{Схемная сложность связности графа}

\( CONN : \{ 0, 1 \}^{\binom{n}{2}} \to \{ 0, 1 \} \)

\( CONN(G) = 1 \) тогда и только тогда, когда \( G \) --- связен.

Какова \( c(CONN) \) ?

\begin{defn}{}{}
    Для неориентированного графа \( g \) его матрица смежности:
    \[
        A_G (i, j) = \begin{cases}
            1, \ \text{\(i, j\) смежны}
            \\
            0, \ \text{иначе}
        \end{cases}
    \]
\end{defn}

\begin{lemma}{}{}
    \( A_G^k (i, j) = \) количество путей из \( i \) в \( j \) длины \( k \).
\end{lemma}

Доказывается индукцией по \( k \) и определением умножения матриц.

Заметим, что если \( G \) связен, то между любыми вершинами \( i \) и \( j \) есть путь длины не более \( n - 1 \).
Иначе говоря, в какой-то из матриц \( A, A^2, \ldots, A^{n - 1} \) для каждой позиции \( (i, j) \) должна быть ненулевая ячейка.

Заменим сложение чисел на \( \lor \), а умножение на \( \land \),
и будем так умножать матрицы (булево умножение).
В силу свойств \( \lor, \land \) такое умножение матриц останется коммутативным, ассоциативным, дистрибутивным.

Получим сложность \( O(n^4) \) на получение матриц \( A_1, \ldots, A^{n - 1} \),
после чего возьмем \( \lor \) по каждой недиагональной позиции по всем недиагональным элементам,
а потом возьмем \( \land \) от полученных значений.
Получили \( O(n^4) \).

Однако если изначально сделать \( A_{ij} = 1 \) (матрица с петлями обозначается \( \tilde{A} \)),
то достаточно проверить только то, что в \( A^{n - 1} \) все элементы равны единице.

А значит нам хватит \( O(n^3 \log n) \) (используем быстрое возведение в степень).

Доказали, что \( c(CONN) = O(n^3 \log n) \)

% \subsection{Глубина булевой схемы}

\begin{defn}{}{}
    Пусть \( x_1, \ldots, x_n, g_1, \ldots, g_s \) --- булева схема,
    Глубина \( g_i \) --- максимальная длина пути от \( x_j \) до \( g_i \)
    (в ориентированном графе булевой схемы).
    Глубина схемы --- максимальная глубина \( g_i \).
\end{defn}

\begin{example}{}{}
    У \( (x_1 \lor (x_2 \lor (x_3 \lor x_4))) \) глубина \( 3 \)

    У \( (x_1 \lor x_2) \lor (x_3 \lor x_4) \) глубина \( 2 \)
\end{example}

% \subsection{Глубина CONN}

Поймем, схему какой глубины мы можем сделать для умножения двух матриц.
Поскольку для каждого элемента нужно ``логическое или'' \( n \) умножений,
можно добиться глубины \( O(\log n) \) с сохранением размера \( O(n^3) \).

Поскольку нам нужно \( O(\log n) \) раз возвести \( \tilde{A_G} \) в квадрат,
а потом взять или всех элементов, получим размер \( O(n^3 \log n) \)
и глубину \( O(\log^2 n) \).

Научимся немного уменьшать размер (но при этом пожертвуем глубиной).
Заметим, что нам достаточно убедиться,
что в \( \tilde{A_G}^{n - 1} \) только первый столбец состоит из единиц.

Для этого нужно найти:
\[
    \tilde{A_G}^{n - 1}
    \cdot
    \begin{pmatrix}
        1
        \\
        0
        \\
        \vdots
        \\
        0
    \end{pmatrix}
\]

Заметим, что если делать умножения не слева направо, а справо налево,
то нам понадобится только \( O(n^3) \) умножений переменных ---
так как каждый раз будем умножать матрицу на столбец.
Однако глубина теперь будет \( O(n \log n) \).

\subsection{Оценки на схемную сложность \( XOR_n \)}

\begin{thrm}{}{}
    \[
        2n - 1 \leq c(XOR_n) \leq 5n - 5 \quad (n \geq 2)
    \]
\end{thrm}

Оценка сверху очевидна, так как уже знаем, что \( c(XOR_2) \leq 5 \).

По индукции докажем, что ни \( XOR_n \), ни \( \lnot XOR_n \) нельзя получить за \( 2n - 2 \)
функциональных элемента.
При этом функциональным элементом будем считать: \( (h_i^a \land h_j^b)^c \)
(\( \lor, \land, \lnot \) выражаются через один такой элемент,
а значит нам хватит таких условий)

\begin{itemize}
    \item
        База \( n = 2 \):

        \( h_1 = (x^a \land y^b)^c \)

        Этого не хватит, так как \( h_1 \) ложен либо на \( 1 \), либо на \( 3 \) наборах.

        Использовать в \( h^2 \) дважды \( h_1 \) не имеет смысла (получим либо \( h_1^a \), либо \( const \)),
        значит \( h_2 = ((x^a \land y^b)^c \land x^d)^e \).
        Переберем случаи, сойдется.
    \item
        Переход:

        Пусть \( XOR_{n + 1} \) можно получить менее чем за \( 2n \) функциональных элемента
        (с \( \lnot XOR_{n + 1} \) поступим аналогично).

        Возьмем булеву схему \( XOR_{n + 1} \).
        Рассмотрим \( g_1 = (x_1^a \land x_2^b)^c \).
        Он подается на вход какому-то другому элементу (иначе его можно выкинуть),
        значит существует \( g_2 = (g_1^{a'} \land h^{b'})^{c'} \).
        Тогда \( g_1 \) станет константой.
        А \( g_2 \) превратится либо в \( h^d \), либо в константу.
        Давайте выкинем \( x_1, g_1, g_2 \).
        Если что-то зависело от \( g_1 \), то оно поведет себя как \( g_2 \),
        а значит можно либо напрямую использовать один из его входов, либо получить константу с помощью
        каких-то двух переменных без использования \( x_1 \).
        С \( g_2 \) аналогично.

        Но тогда мы научились считать \( XOR_n, \lnot XOR_n \) за \( 2n - 3 \) --- противоречие.
\end{itemize}

\begin{exercise}{}{}
    \[
        c(XOR_N) \geq 3n - 3
    \]
\end{exercise}

\subsection{Задачи выполнимости}

Выполнимость --- существование \( x \in \{ 0, 1 \}^n : f(x) = 1 \)

\begin{enumerate}
    \item
        Дана схема \( C(x) \) от  \( n \) переменных, нужно проверить, выполнима ли она.
    \item
        Дана КНФ \( K(x) \), нужно проверить, выполнима ли она.
    \item
        Дана 3-КНФ (каждый дизъюнкт из не более чем трех переменных) \( K(x) \), нужно проверить, выполнима ли она 
\end{enumerate}

Заметим, что сведения \( (3) \to (2) \to (1) \) очевидны (так как по сути это частные случаи).

\begin{thrm}{}{}
    Задача выполнимости булевой схемы эффективно сводится к задаче выполнимости 3-КНФ.
\end{thrm}

Пусть есть схема \( x_1, \ldots, x_n, g_1, \ldots, g_s \).

Хотим получить 3-КНФ \( K(y) \) размера \( K = poly(s) \).

\begin{itemize}
    \item
        \(
            g_i = \lnot h_j
            \Leftrightarrow
            (g_i \lor h_j) \land (\lnot g_i \lor \lnot h_j)
        \)
    \item
        \(
            g_i = h_k \land h_l
            \Leftrightarrow
            (h_k \lor \lnot g_i) \land (h_l \lor \lnot g_i) \land (\lnot h_k \lor \lnot h_l \lor g_i)
        \)
    \item
        \(
            g_i = h_k \lor h_l
            \Leftrightarrow
            (\lnot h_k \lor g_i) \land (\lnot h_l \lor g_i) \land (h_k \lor h_l \lor \lnot g_i)
        \)
\end{itemize}

Пройдемся по схеме и запишем описанные выше 3-КНФ.
Если полученная 3-КНФ разрешима, то явно существует \( x : C(x) = 1 \).
В обратную сторону тоже тривиально.
При этом размер 3-КНФ \( O(s) \),
так как все операции из схемы мы делаем за \( O(1) \) дизъюнктов.
